\documentclass[11pt]{article}

\usepackage[utf8]{inputenc}
\usepackage[T1]{fontenc}
\usepackage{lmodern}
\usepackage{geometry}
\usepackage{hyperref}
\usepackage{xcolor}
\usepackage{longtable}
\usepackage{array}
\usepackage{listings}
\usepackage{courier}
\geometry{margin=1in}

\hypersetup{
    colorlinks=true,
    linkcolor=blue,
    urlcolor=blue
}

\lstset{
    basicstyle=\ttfamily\small,
    breaklines=true,
    frame=single,
    columns=fullflexible
}

\title{Technical Brief:\\The LogicScript Optimizing Compiler}
\author{UFUG2106 Project Team}
\date{\today}

\begin{document}

\maketitle

\section{Project Overview}
This project implements a complete four-phase compiler pipeline for the LogicScript micro-language described in \texttt{UFUG2106-Project.pdf}. The compiler reads a human-readable text program, validates it, optimizes its logical expressions, proves that the optimizations are safe, and then executes the optimized program. The final output is exported as a JSON trace so that the grading script can inspect every successful phase and, if necessary, the exact phase and line where compilation failed.

The implementation follows the assignment structure directly. Each major compiler phase is isolated in its own Python file:
\begin{itemize}
    \item \texttt{phase\_1\_lexer.py}
    \item \texttt{phase\_2\_parser.py}
    \item \texttt{phase\_3\_optimizer.py}
    \item \texttt{phase\_4\_execution.py}
\end{itemize}
The coordinator file \texttt{main.py} connects these modules into one full pipeline, while the root-level \texttt{logic\_compiler.py} provides the command-line interface required by the assignment.

\section{Team Organization}
The project was organized by compiler phase ownership, while still keeping the full pipeline jointly reviewed and integrated. Since each phase depends on the previous one, dividing work only by file was not enough; the team also had to maintain a shared understanding of the token format, AST structure, optimization output, and final JSON layout.

\subsection{Bi Li}
Bi Li was mainly responsible for the optimizer part of the compiler, that is the file \texttt{LogicScript/phase\_3\_optimizer.py}. He implemented and refined the logical simplification rules, including the associative normalization extension for nested \texttt{AND} and \texttt{OR} expressions. He also wrote a major part of the report and reviewed the overall consistency of the compiler design.

\subsection{Olzhas Tilektes}
Olzhas was mainly responsible for the lexer and parser parts of the compiler, that is the files \texttt{LogicScript/phase\_1\_lexer.py} and \texttt{LogicScript/phase\_2\_parser.py}. He worked on the tokenization rules, the recursive-descent parsing structure, and the overall integration of the early compiler pipeline. He also coordinated the project structure and edited files during integration.

\subsection{Kenn Febrain Susmawanto}
Kenn was mainly responsible for the verification and execution part of the compiler, that is the file \texttt{LogicScript/phase\_4\_execution.py}. He worked on truth-table verification, expression evaluation, execution with the state dictionary, and the final JSON trace generation. He also contributed to the project documentation, especially the README and submission materials.

Even though the work was divided by phase, the final compiler required repeated joint checking. The parser output had to match what the optimizer expects, the optimizer had to preserve a structure that the executor can evaluate, and the final JSON format had to remain consistent with the grading requirements. Therefore the contributions were phase-based, but the final result was achieved collaboratively.

\section{System Design}
\subsection{Overall Pipeline}
The compiler processes the input program in four sequential stages:
\begin{enumerate}
    \item \textbf{Lexical analysis}: convert each source line into canonical tokens.
    \item \textbf{Parsing}: validate the grammar and build nested AST lists.
    \item \textbf{Optimization}: simplify AST expressions using propositional logic laws.
    \item \textbf{Verification and execution}: prove changed expressions equivalent by truth table, then execute the optimized program.
\end{enumerate}

The modular coordinator is \texttt{LogicScript/main.py}. If we treat the four phases as black boxes, then \texttt{main.py} is the file that connects them together in the correct order and collects their outputs into one final dictionary. It imports:
\begin{itemize}
    \item \texttt{run\_lexer} from phase 1
    \item \texttt{run\_parser} from phase 2
    \item \texttt{run\_optimizer} from phase 3
    \item \texttt{run\_execution} from phase 4
    \item \texttt{CompilerError} for graceful error handling
\end{itemize}

The central function is \texttt{compile\_source\_lines(source\_lines)}. It:
\begin{enumerate}
    \item creates an empty result dictionary
    \item runs phase 1 and stores \texttt{result["phase\_1\_lexer"]}
    \item runs phase 2 and stores \texttt{result["phase\_2\_parser"]}
    \item runs phase 3 and stores \texttt{result["phase\_3\_optimizer"]}
    \item runs phase 4 and stores \texttt{result["phase\_4\_execution"]}
    \item returns the full result dictionary
\end{enumerate}

One important part of \texttt{main.py} is the error handling. The whole pipeline is wrapped in a \texttt{try/except CompilerError} block. If any phase fails, the compiler does not crash with a Python traceback. Instead, it returns the outputs of the successful earlier phases together with:
\begin{lstlisting}[language=Python]
"error": {
    "phase": error.phase,
    "line": error.line,
}
\end{lstlisting}
This behavior matches the assignment requirement that the compiler halt gracefully and report the failing phase and line number.

\subsection{Phase 1: Lexer}
The file \texttt{LogicScript/phase\_1\_lexer.py} handles the lexical part of the compilation. In other words, lexer reads the input plain text file and splits it into pieces. The task of the lexer is to create a dictionary of tokens which is later used for parsing in phase 2.

For example, if an \texttt{input.txt} file contains:
\begin{quote}
\texttt{let p = T}
\end{quote}
lexer returns tokens of the form:
\begin{quote}
\texttt{\{"line": 1, "tokens": ["LET", "VAR\_P", "EQ", "TRUE"]\}}
\end{quote}

Lexer also catches lexical errors such as unconventional expressions:
\begin{quote}
\texttt{le tp = T}
\end{quote}
This causes a \texttt{CompilerError} because the separated piece \texttt{le} does not belong to the LogicScript alphabet.

More precisely, lexer reads the file line-by-line and separates each line into pieces. A piece must belong to one of five lexical categories:
\begin{enumerate}
    \item booleans: \texttt{T}, \texttt{F}
    \item variables: any single lowercase letter
    \item symbols: \texttt{(}, \texttt{=}, \texttt{)}
    \item operators: \texttt{NOT}, \texttt{AND}, \texttt{OR}, \texttt{IMPLIES}
    \item keywords: \texttt{let}, \texttt{if}, \texttt{then}, \texttt{print}
\end{enumerate}

The file uses dictionaries to classify each reserved token:
\begin{lstlisting}[language=Python]
KEYWORDS = {"let": "LET", "if": "IF", "then": "THEN", "print": "PRINT"}
BOOLEANS = {"T": "TRUE", "F": "FALSE"}
OPERATORS = {"AND": "AND", "OR": "OR", "NOT": "NOT", "IMPLIES": "IMPLIES"}
SYMBOLS = {"=": "EQ", "(": "L_PAREN", ")": "R_PAREN"}
\end{lstlisting}

It then uses a regular expression scanner:
\begin{lstlisting}
r"\s+|[()]|=|[A-Za-z]+|."
\end{lstlisting}
This regex separates whitespace, parentheses, equals signs, letter-runs, and any leftover single invalid character. The final \texttt{.} in the regex is important because it prevents illegal symbols from being silently skipped.

The three main lexer functions are:
\begin{itemize}
    \item \texttt{classify\_token}: classify one piece
    \item \texttt{tokenize\_line}: tokenize one line
    \item \texttt{run\_lexer}: tokenize the whole source file
\end{itemize}

\subsection{Phase 2: Parser}
The file \texttt{LogicScript/phase\_2\_parser.py} handles parsing after lexical analysis is completed. Parser reads the standardized tokens that were produced in phase 1 and checks whether they follow the recursive grammar of LogicScript. The task of the parser is to create a dictionary of ASTs which is later used for optimization in phase 3.

For example, if lexer returned:
\begin{quote}
\texttt{\{"line": 1, "tokens": ["LET", "VAR\_P", "EQ", "TRUE"]\}}
\end{quote}
parser produces:
\begin{quote}
\texttt{\{"line": 1, "ast": ["LET", "VAR\_P", "TRUE"]\}}
\end{quote}

Parser also catches grammatical errors that are not lexical errors. For example:
\begin{quote}
\texttt{if (p AND ) then print p}
\end{quote}
is lexically valid but grammatically invalid, because after the binary operator \texttt{AND} there is no valid right-hand expression.

Parser builds a nested list structure called the AST. In this representation, parentheses and \texttt{EQ} are removed, while the logical grouping is preserved directly by nested Python lists. The recursive grammar rules are:
\begin{enumerate}
    \item base expressions: \texttt{TRUE}, \texttt{FALSE}, variables
    \item unary expressions: \texttt{(NOT E)}
    \item binary expressions: \texttt{(E1 AND E2)}, \texttt{(E1 OR E2)}, \texttt{(E1 IMPLIES E2)}
    \item statements: \texttt{let}, \texttt{if ... then ...}, \texttt{print}
\end{enumerate}

The parser is implemented as a recursive-descent parser. It uses two small helper sets:
\begin{lstlisting}[language=Python]
BOOLEAN_LITERALS = {"TRUE", "FALSE"}
BINARY_OPERATORS = {"AND", "OR", "IMPLIES"}
\end{lstlisting}
and a helper function:
\begin{lstlisting}[language=Python]
def is_variable(token):
    return isinstance(token, str) and token.startswith("VAR_")
\end{lstlisting}

The core data structure of phase 2 is the \texttt{TokenStream} class, which stores the token list, the source line number, and the current parser position. This class provides:
\begin{itemize}
    \item \texttt{current()} for lookahead
    \item \texttt{consume()} for controlled token consumption
    \item \texttt{consume\_variable()} for variable-only positions
\end{itemize}

The main recursive function is \texttt{parse\_expression(stream)}. It checks whether the current token starts a base expression or a parenthesized recursive expression, and then returns the corresponding AST. Statement parsing is handled by a dispatcher that maps the first token of the line to one of:
\begin{itemize}
    \item \texttt{parse\_let\_statement}
    \item \texttt{parse\_if\_statement}
    \item \texttt{parse\_print\_statement}
\end{itemize}

At the line level, parser uses \texttt{parse\_line(tokens, line\_number)} to ensure that one token list forms exactly one complete statement and contains no trailing unused tokens.

\subsection{Phase 3: Optimizer}
The file \texttt{LogicScript/phase\_3\_optimizer.py} handles optimization. Optimizer reads the ASTs produced in phase 2 and simplifies the logical expressions inside them. The task of the optimizer is to create optimized ASTs and also collect proof inputs for phase 4.

For example, if parser returned:
\begin{quote}
\texttt{\{"line": 1, "ast": ["LET", "VAR\_X", ["OR", "VAR\_P", "TRUE"]]\}}
\end{quote}
optimizer returns:
\begin{quote}
\texttt{\{"line": 1, "ast": ["LET", "VAR\_X", "TRUE"]\}}
\end{quote}

Optimizer is a static pass. It does not execute the program and does not look at runtime values stored in the state dictionary. Instead, it recursively traverses each expression subtree and applies logical equivalence rules. It returns two outputs:
\begin{lstlisting}[language=Python]
phase_3_output = [
    {"line": line_number, "ast": optimized_ast}
]

verification_pairs = [
    (line_number, original_expr, optimized_expr)
]
\end{lstlisting}

Only expression parts are optimized. Statement shapes are preserved:
\begin{itemize}
    \item \texttt{LET}: optimize only the right-hand expression
    \item \texttt{IF}: optimize the condition and nested statement
    \item \texttt{PRINT}: unchanged
\end{itemize}

The optimizer implements propositional logic laws such as:
\begin{itemize}
    \item constant folding
    \item identity
    \item double negation
    \item De Morgan's laws
    \item complement
    \item absorption
    \item implication rewriting
\end{itemize}

It is a rule-based optimizer, not a global Boolean minimizer. Therefore it does not guarantee the mathematically simplest possible equivalent expression in every case.

The file uses several helper predicates such as:
\begin{lstlisting}[language=Python]
def is_true(expr):
    return expr == "TRUE"

def is_false(expr):
    return expr == "FALSE"

def is_not(expr):
    return isinstance(expr, list) and len(expr) == 2 and expr[0] == "NOT"
\end{lstlisting}
and a more general binary matcher:
\begin{lstlisting}[language=Python]
def is_binary(expr, operator=None):
    if not (isinstance(expr, list) and len(expr) == 3):
        return False
    if operator is None:
        return True
    return expr[0] == operator
\end{lstlisting}

Each logical law is implemented as an ordinary Python function. For example:
\begin{lstlisting}[language=Python]
def rule_and_identity(left, right):
    if is_true(left):
        return right
    if is_true(right):
        return left
    return None
\end{lstlisting}
If a rule matches, it returns the simplified expression. Otherwise it returns \texttt{None}. Rules are tried in order using:
\begin{lstlisting}[language=Python]
def apply_rule_list(*args, rules):
    for rule in rules:
        result = rule(*args)
        if result is not None:
            return result
    return None
\end{lstlisting}

The core recursive function is \texttt{optimize\_expression(expr)}, which works bottom-up: optimize the children first, then optimize the current node. This is why many multi-step simplifications can still occur even without a separate global fixed-point pass.

One contribution added by Bi Li is the associative normalization for nested \texttt{AND} and \texttt{OR} chains. The helper functions \texttt{collect\_associative\_operands(...)} and \texttt{build\_associative\_expression(...)} flatten and rebuild equivalent chains into one canonical shape, which helps later laws such as absorption trigger more often.

Finally, optimizer collects \texttt{(line\_number, original\_expr, optimized\_expr)} pairs for every changed expression. These pairs are then checked by phase 4.

\subsection{Phase 4: Verification and Execution}
The file \texttt{LogicScript/phase\_4\_execution.py} handles the final part of the compilation. Phase 4 reads the optimized ASTs produced in phase 3, proves that every optimization is logically safe, and then executes the optimized program. The task of this phase is to create the execution trace that later appears in the final JSON output.

Phase 4 has two subparts:
\begin{enumerate}
    \item verification
    \item execution
\end{enumerate}

Verification is the safety net of the compiler. Whenever optimizer changed an expression, phase 4 generates a full truth table for the original expression and the optimized expression. If the two output columns differ, phase 4 raises an error immediately.

Execution then runs the optimized statements line-by-line using a state dictionary. This dictionary stores the current boolean value of each variable. Printed values are not written directly to standard output; instead they are stored in the structured JSON trace required by the assignment.

The final return value has the form:
\begin{lstlisting}[language=Python]
{
    "verifications": [...],
    "final_state_dictionary": {...},
    "printed_output": [...]
}
\end{lstlisting}

The file uses two main ideas:
\begin{enumerate}
    \item recursive evaluation of AST expressions
    \item recursive generation of all truth-table assignments
\end{enumerate}

Binary truth functions are stored in:
\begin{lstlisting}[language=Python]
EXPRESSION_EVALUATORS = {
    "AND": lambda left, right: "TRUE" if left == "TRUE" and right == "TRUE" else "FALSE",
    "OR": lambda left, right: "TRUE" if left == "TRUE" or right == "TRUE" else "FALSE",
    "IMPLIES": lambda left, right: "FALSE" if left == "TRUE" and right == "FALSE" else "TRUE",
}
\end{lstlisting}

To construct a truth table, phase 4 first collects the variables that appear in the expression tree:
\begin{lstlisting}[language=Python]
def collect_variables(expr):
    if isinstance(expr, str):
        return {expr} if is_variable(expr) else set()
    if len(expr) == 2:
        return collect_variables(expr[1])
    return collect_variables(expr[1]) | collect_variables(expr[2])
\end{lstlisting}
It then generates all truth assignments recursively with \texttt{generate\_truth\_assignments(variables)}. If there are \(n\) variables, then \(2^n\) assignments are generated, which is exactly what a full truth table requires.

The heart of both verification and execution is \texttt{eval\_expression(expr, state, line\_number)}. This function:
\begin{itemize}
    \item returns literals directly
    \item looks up variables in the state dictionary
    \item recursively evaluates \texttt{NOT}
    \item recursively evaluates binary expressions and applies the correct truth function
\end{itemize}

After all verification pairs are proven equivalent, execution proceeds using three statement handlers:
\begin{itemize}
    \item \texttt{execute\_let\_statement}
    \item \texttt{execute\_print\_statement}
    \item \texttt{execute\_if\_statement}
\end{itemize}

Assignment updates the state dictionary, print appends a structured output record, and \texttt{IF} evaluates the condition and executes its nested statement only if the condition is \texttt{"TRUE"}.

\section{Testing Suite}
The file \texttt{LogicScript/tests.py} contains the local testing suite of the project. It is intentionally simple and does not use an external testing framework. Instead it uses ordinary Python functions together with \texttt{assert} statements. This is sufficient for this project because the goal is to verify that the whole compiler pipeline behaves correctly on several representative inputs.

The structure of the file is:
\begin{enumerate}
    \item import \texttt{compile\_source\_lines} from \texttt{main.py}
    \item construct small LogicScript programs directly as Python lists of strings
    \item run the full compiler in memory
    \item compare the actual result with the expected result
\end{enumerate}

\subsection{Successful Compilation Case}
The strongest end-to-end test is \texttt{test\_valid\_program()}, which uses the same successful example as the assignment:
\begin{quote}
\texttt{let p = T}\\
\texttt{let q = F}\\
\texttt{let r = (NOT ((NOT p) AND q))}\\
\texttt{if r then print p}
\end{quote}
This test checks the exact outputs of all four phases:
\begin{itemize}
    \item token output of phase 1
    \item AST output of phase 2
    \item optimized AST output of phase 3
    \item verification, final state dictionary, and printed output of phase 4
\end{itemize}

\subsection{Syntax Error Case}
The second test, \texttt{test\_syntax\_error()}, checks whether the compiler handles a grammatical error properly:
\begin{quote}
\texttt{let p = T}\\
\texttt{if (p AND ) then print p}
\end{quote}
This line is lexically valid, so phase 1 should still succeed. However parser should detect that the binary operator \texttt{AND} has no valid right-hand expression after it. The expected result is:
\begin{lstlisting}[language=Python]
{"error": {"phase": "phase_2_parser", "line": 2}}
\end{lstlisting}

\subsection{Lexical Error Case}
The third test, \texttt{test\_lexical\_error()}, checks whether phase 1 catches an invalid raw symbol:
\begin{quote}
\texttt{let p = T}\\
\texttt{let q = @}
\end{quote}
Here the symbol \texttt{@} is not part of the LogicScript alphabet, so phase 1 must halt immediately and report:
\begin{lstlisting}[language=Python]
{"error": {"phase": "phase_1_lexer", "line": 2}}
\end{lstlisting}

\subsection{Optimization Verification Case}
The fourth test, \texttt{test\_constant\_folding\_verification()}, focuses on a simple optimization:
\begin{quote}
\texttt{let p = T}\\
\texttt{let x = (p OR T)}\\
\texttt{print x}
\end{quote}
The expected behavior is:
\begin{itemize}
    \item phase 3 simplifies \texttt{(p OR T)} into \texttt{TRUE}
    \item phase 4 records a verification object proving that the original and optimized expressions are equivalent
    \item the printed result is \texttt{TRUE}
\end{itemize}

\subsection{What the Tests Show}
The tests can be run from the \texttt{LogicScript} directory with:
\begin{lstlisting}
python tests.py
\end{lstlisting}
These tests are useful because they check:
\begin{itemize}
    \item successful end-to-end compilation
    \item lexical error handling
    \item syntax error handling
    \item one optimization together with truth-table verification
\end{itemize}
At the same time, the current test suite is still limited. It does not yet cover every possible edge case, such as more complex optimizer interactions, additional undefined-variable execution failures, larger families of parser corner cases, and more cases for the associative optimization added in phase 3.

\section{Limitations}
The language and compiler are intentionally simple, which makes them clear and suitable for the project requirements. However, this simplicity also creates several important limitations.

\subsection{Restricted Variable System}
At the moment variables can only be single lowercase letters. This means the language supports at most 26 direct variable names such as \texttt{a}, \texttt{b}, \texttt{p}, and \texttt{x}. While this follows the project specification, it is restrictive from a usability point of view. A natural future improvement would be to let lexer accept longer identifier names such as words or alphanumeric variable names.

\subsection{Verbose Parenthesized Syntax}
LogicScript requires explicit parentheses for recursive expressions. Even a moderately simple formula may look like:
\begin{quote}
\texttt{let x = (((NOT a) AND b) OR c)}
\end{quote}
This is correct and unambiguous, but it is not very human-friendly. A future improvement could be to add precedence rules so that \texttt{NOT} binds tighter than \texttt{AND}, and \texttt{AND} binds tighter than \texttt{OR}. This would allow less cluttered expressions, but it would also require a significantly more complex parser.

\subsection{Line-by-Line Parsing}
The current parser assumes one full statement per line, which means formatting flexibility is low. For example, a user cannot write:
\begin{quote}
\texttt{let x = a AND}\\
\texttt{\ \ \ \ \ \ \ \ (}\\
\texttt{\ \ \ \ \ \ \ \ \ \ \ \ b}\\
\texttt{\ \ \ \ \ \ \ \ )}
\end{quote}
The compiler would not treat this as one statement. This is not a problem for the required project grammar, but it limits readability and scalability. Supporting multi-line statements would require redesigning how the compiler groups lines before parsing.

\subsection{Rule-Based Optimization}
The optimizer is rule-based, which means it applies only the specific logical equivalence rules that were explicitly implemented. This is enough for many useful simplifications, but it does not guarantee the mathematically simplest equivalent expression in every case.

For example, consider:
\begin{quote}
\texttt{let x = ((NOT p) AND ((NOT q) AND (NOT r)))}\\
\texttt{let y = (p OR (q AND r))}\\
\texttt{let z = (x OR y)}
\end{quote}
The optimizer can simplify many local patterns, but it may still leave an expression in a non-minimal form if the simplification would require more global Boolean reasoning. In principle, stronger minimization methods such as Karnaugh maps or Quine-McCluskey tables could reduce some expressions further. However, these methods are considerably more complex and go beyond the intended scope of the current project.

Therefore, the current compiler should be understood as a correct and structured educational compiler, not as a fully general Boolean minimization system. Its main strength is clarity, modularity, and alignment with the course concepts of recursion, propositional logic, and state relations.

\section{Conclusion}
The LogicScript compiler demonstrates how ideas from discrete mathematics can be turned into a working software tool. Recursive definitions drive the parser, propositional equivalence laws guide optimization, and state dictionaries support both execution and verification. The final system is small, but it reflects the main structure of a real compiler: representation, transformation, proof of correctness, and controlled execution.

\end{document}
